\documentclass[10pt]{resume}
\usepackage[left=0.4in,top=0.3in,right=0.4in,bottom=0.3in]{geometry}
\usepackage[colorlinks=true,urlcolor=blue]{hyperref}
\usepackage{enumitem}

\name{LOÏC ARON MBASSI EWOLO}
\address{Alternance Ingénierie Système : Véhicule Autonome / GNSS (dès Septembre 2026)}
\address{
Cergy, France \\
\href{mailto:loicaron.mbassiewolo@ensea.fr}{loicaron.mbassiewolo@ensea.fr} \,|\,
\href{tel:+33759523398}{+33 7 59 52 33 98} \\
\href{https://www.linkedin.com/in/loic-mbassi/}{LinkedIn} \,|\,
\href{https://github.com/Nameless0l}{GitHub}
}

\begin{document}

%--------------------------------------------------
% RÉSUMÉ PROFESSIONNEL / PROFIL
%--------------------------------------------------
\begin{rSection}{Résumé Professionnel}
Élève-ingénieur en double diplôme ENSEA/Polytechnique Yaoundé, spécialisé en systèmes embarqués et navigation autonome. Expérience en développement C/C++, spécifications fonctionnelles et simulation. Passionné par les systèmes de mobilité innovants et la géolocalisation GNSS.
\end{rSection}

%--------------------------------------------------
\begin{rSection}{Éducation}
{\bf ENSEA - École Nationale Supérieure de l'Électronique} \hfill {\em 2025--2027}\\
\textbf{Double diplôme d'ingénieur -- Systèmes Embarqués \& IA}

{\bf École Nationale Supérieure Polytechnique de Yaoundé (ENSPY)} \hfill {\em 2021--2025}\\
\textbf{Diplôme d'ingénieur de conception en informatique} (Niveau Master 2)
\end{rSection}

%--------------------------------------------------
\begin{rSection}{Compétences Techniques}
\begin{tabular}{ @{} >{\bfseries}l @{\hspace{6ex}} l }
Langages & C, C++, C\#, Python, Java \\
Embarqué \& Temps réel & STM32, Arduino, Nvidia Jetson, Linux embarqué, IPC, Signaux \\
Navigation \& Capteurs & GNSS, Lidar 3D, SLAM, Fusion capteurs, Caméra de profondeur \\
Outils & Git/GitLab, Docker, VS Code, Modélisation 3D (Onshape) \\
Langues & Français (Natif), Anglais (Professionnel/Technique)
\end{tabular}
\end{rSection}

%--------------------------------------------------
\begin{rSection}{Projets Académiques}

{\bf Upgrade Jetson -- Challenge CoHoMa (Robotique Défense)} \hfill {\em 2025}\\
Développement d'un robot autonome militaire en C/C++ dans le cadre d'un challenge avec l'Armée Française. Intégration d'un Lidar 3D et d'une caméra de profondeur pour la cartographie. Implémentation des algorithmes SLAM et fusion de capteurs sur Nvidia Jetson.

{\bf PROJET-TAXI -- Simulation Multiprocessus (C/Linux)} \hfill {\em 2024}\\
Conception d'un simulateur de flotte de taxis en C sous Linux. Utilisation de la mémoire partagée (SHM) et des signaux POSIX pour la communication inter-processus. Rédaction des spécifications fonctionnelles et définition des interfaces.

{\bf Fault Detection \& Control (Drone) -- MATLAB/Simulink} \hfill {\em 2025}\\
Modélisation dynamique d'un drone en conditions nominales et dégradées. Conception d'un schéma de détection de défauts basé sur des observateurs. Mise en place d'une stratégie de contrôle tolérant aux pannes.

\end{rSection}

%--------------------------------------------------
\begin{rSection}{Expérience Professionnelle}
{\bf Assistant de Recherche -- MILA (Institut Québécois d'IA)} \hfill {\em Juin 2025 -- Sept 2025}\\
Étude du phénomène de "Grokking" (généralisation tardive des réseaux de neurones). Reproduction des expériences état de l'art et analyse de la convergence mathématique en Python/PyTorch.

{\bf Stage Recherche IoT Lab -- TinyML, ENSPY} \hfill {\em Oct 2023 -- Jan 2024}\\
Optimisation de modèles de Machine Learning pour le déploiement sur des plateformes à ressources limitées. Développement en C sur Arduino et Raspberry Pi avec des contraintes temps réel.

{\bf Développeur Backend -- CSC (Fintech), Yaoundé} \hfill {\em Oct 2024 -- Jan 2025}\\
Conception du backend d'une plateforme de transactions financières sécurisée. Mise en place d'une architecture scalable avec Docker et MySQL.
\end{rSection}

%--------------------------------------------------
\begin{rSection}{Distinctions et Engagements}

\textbf{1\textsuperscript{er} Prix CITS National Hackathon 2024} -- Optimisation routes (algorithme voyageur de commerce). \\
\textbf{1\textsuperscript{er} Prix IndabaX Africa 2025} -- IA Santé, déploiement API Flask/Streamlit. \\
\textbf{Lead AI Cell, Polytechnique Yaoundé} -- +50\% membres, collaboration Google DeepMind.
\end{rSection}

\end{document}
